\section{Method}
\label{sec:method}

\subsection{The streaming system}

The system is a three-thread proactive pipeline over a single shared KV cache. An
\emph{ingester} decodes video at 1\,fps and audio in 2\,s chunks and appends both
to the cache. A \emph{controller} wakes on every tick, and a \emph{writer} decodes
an utterance whenever the controller decides to speak. The backbone is
Qwen2.5-Omni-7B; the same harness previously ran a vision-only Qwen3-VL-8B, and
the pipeline is unchanged apart from the backbone and the audio path.

At each tick the controller poses a boolean question to the model and reads
$p_{\mathrm{hit}} = P(\texttt{true})$ directly off the logits at the answer slot.
This costs one forward pass and \emph{zero} decode steps, which is what makes a
per-tick gate affordable: the expensive answer decode happens only when the gate
fires. The consequence for cost is that a low threshold is substantially more
expensive than a high one, because it admits more decodes --- measured at
$\mathrm{xRT}=2.72$ at threshold $0.05$ against a quiet-path cost of
$0.3$\,s per tick with zero generated tokens.

\subsection{The gate is three coupled knobs, not one}

The decision to speak is governed by three parameters that cannot be fitted
independently:

\begin{itemize}
\item \textbf{\texttt{hit\_threshold}} --- the level $p_{\mathrm{hit}}$ must reach.
\item \textbf{\texttt{gate\_mode}} --- \emph{level} (fire whenever above) or
      \emph{edge} (fire only on a rising crossing).
\item \textbf{\texttt{refractory\_s}} --- the minimum silence after an emission.
\end{itemize}

The coupling is not incidental. Under \emph{edge} mode with a long refractory, a
single spurious early crossing can consume a video's entire emission budget, so
the cost of a mis-set threshold scales with the refractory. A threshold reported
without its mode and refractory does not reproduce.

\subsection{Two defects in the shipped gate, and their repair}

Auditing the shipped configuration before fitting anything revealed that the
quantity we intended to fit was not the quantity the system used.

\paragraph{The threshold was inert.} Under the shipped
\texttt{gate\_strategy="hysteresis"}, the fire test was
$p_{\mathrm{hit}} \ge \texttt{gate\_high\_thr}$, a fixed global $0.5$;
\texttt{hit\_threshold} only decided whether an answer was subsequently decoded,
and \texttt{gate\_mode} and \texttt{refractory\_s} were read nowhere at all.
Across a completed 2{,}700-sample run (581{,}403 ticks, 50{,}433 fires) not one
fire occurred below $p_{\mathrm{hit}}=0.500$. Every task configured at or under
$0.5$ was therefore running at $0.5$. We add a \texttt{"fitted"} strategy in which
all three knobs are live, and report the shipped behaviour honestly as
\emph{a fixed global $0.5$} rather than as the per-task table that was nominally
configured.

\paragraph{The first tick was a spurious rising edge.} $p_{\mathrm{hit}}$ is high
on the first tick of every video: the model answers the boolean from the prompt
before it has seen anything ($\ge 0.05$ for 98--100\,\% of samples in every task,
median $0.34$--$0.64$). Seeding the edge detector with
$\texttt{prev\_above}=\textsc{false}$ made that tick a rising edge. For a task with
a long refractory this consumed the video's only emission: every
\texttt{instant\_event\_alert} sample emitted once, at $1.0$\,s, against ground
truth at 26--98\,s. The fix is definitional rather than a tuned guard --- an edge
requires a predecessor, so $\texttt{prev\_above}$ starts undefined and the first
tick can never be a crossing. Offline best time-F1 rose $0.0535 \to 0.0900$ on
\texttt{snapshot\_counting} and $0.0636 \to 0.0751$ on
\texttt{instant\_event\_alert}, while short-refractory tasks moved by $\pm 0.001$.
That the size of the effect tracks the refractory is the practical demonstration
that the three knobs are one object.
